% Chapter 1: Smooth Manifolds
% Lee, *Introduction to Smooth Manifolds* (2nd ed., GTM 218).
% Generated from the section extraction; nodes carry only Lean that is
% verified sorry-free. See ../../UPSTREAM_LEAN_AUDIT.md.


\section{Topological Manifolds}

\begin{definition}
\label{def:1-extra-2}
\lean{OpenPartialHomeomorph}
\mathlibok
Let $M$ be a topological $n$ -manifold. A coordinate chart (or just a chart) on $M$ is a pair $\left( {U,\varphi }\right)$ , where $U$ is an open subset of $M$ and $\varphi  : U \rightarrow  \widehat{U}$ is a homeomorphism from $U$ to an open subset $\widehat{U} = \varphi \left( U\right)  \subseteq  {\mathbb{R}}^{n}$ (Fig. 1.2). By definition of a topological manifold, each point $p \in  M$ is contained in the domain of some chart $\left( {U,\varphi }\right)$ . If $\varphi \left( p\right)  = 0$ , we say that the chart is centered at $p$ . If $\left( {U,\varphi }\right)$ is any chart whose domain contains $p$ , it is easy to obtain a new chart centered at $p$ by subtracting the constant vector $\varphi \left( p\right)$ .

Given a chart $\left( {U,\varphi }\right)$ , we call the set $U$ a coordinate domain, or a coordinate neighborhood of each of its points. If, in addition, $\varphi \left( U\right)$ is an open ball in ${\mathbb{R}}^{n}$ , then $U$ is called a coordinate ball; if $\varphi \left( U\right)$ is an open cube, $U$ is a coordinate cube. The map $\varphi$ is called a (local) coordinate map, and the component functions $\left( {{x}^{1},\ldots ,{x}^{n}}\right)$ of $\varphi$ , defined by $\varphi \left( p\right)  = \left( {{x}^{1}\left( p\right) ,\ldots ,{x}^{n}\left( p\right) }\right)$ , are called local coordinates on $U$ . We sometimes write things such as " $\left( {U,\varphi }\right)$ is a chart containing $p$ " as shorthand for " $\left( {U,\varphi }\right)$ is a chart whose domain $U$ contains $p$ ." If we wish to emphasize the coordinate functions $\left( {{x}^{1},\ldots ,{x}^{n}}\right)$ instead of the coordinate map $\varphi$ , we sometimes denote the chart by $\left( {U,\left( {{x}^{1},\ldots ,{x}^{n}}\right) }\right)$ or $\left( {U,\left( {x}^{i}\right) }\right)$ .
\end{definition}

\begin{example}[Graphs of Continuous Functions]
\label{ex:1-3}
\lean{graph_coordinates}
\leanok
\dcref{ch1:1.3}
Let $U \subseteq  {\mathbb{R}}^{n}$ be an open subset, and let $f : U \rightarrow  {\mathbb{R}}^{k}$ be a continuous function. The graph of $f$ is the subset of ${\mathbb{R}}^{n} \times  {\mathbb{R}}^{k}$ defined by

\[
\Gamma \left( f\right)  = \left\{  {\left( {x, y}\right)  \in  {\mathbb{R}}^{n} \times  {\mathbb{R}}^{k} : x \in  U\text{ and }y = f\left( x\right) }\right\}  ,
\]

with the subspace topology. Let ${\pi }_{1} : {\mathbb{R}}^{n} \times  {\mathbb{R}}^{k} \rightarrow  {\mathbb{R}}^{n}$ denote the projection onto the first factor, and let $\varphi  : \Gamma \left( f\right)  \rightarrow  U$ be the restriction of ${\pi }_{1}$ to $\Gamma \left( f\right)$ :

\[
\varphi \left( {x, y}\right)  = x,\;\left( {x, y}\right)  \in  \Gamma \left( f\right) .
\]

Because $\varphi$ is the restriction of a continuous map, it is continuous; and it is a homeomorphism because it has a continuous inverse given by ${\varphi }^{-1}\left( x\right)  = \left( {x, f\left( x\right) }\right)$ . Thus $\Gamma \left( f\right)$ is a topological manifold of dimension $n$ . In fact, $\Gamma \left( f\right)$ is homeomorphic to $U$ itself, and $\left( {\Gamma \left( f\right) ,\varphi }\right)$ is a global coordinate chart, called graph coordinates. The same observation applies to any subset of ${\mathbb{R}}^{n + k}$ defined by setting any $k$ of the coordinates (not necessarily the last $k$ ) equal to some continuous function of the other $n$ , which are restricted to lie in an open subset of ${\mathbb{R}}^{n}$ .
\end{example}

\begin{example}[Projective Spaces]
\label{ex:1-5}
\lean{real_projective_space_has_standard_chart}
\leanok
\dcref{ch1:1.5}
The $n$ -dimensional real projective space, denoted by ${\mathbb{{RP}}}^{n}$ (or sometimes just ${\mathbb{P}}^{n}$ ), is defined as the set of 1-dimensional linear subspaces of ${\mathbb{R}}^{n + 1}$ , with the quotient topology determined by the natural map $\pi  : {\mathbb{R}}^{n + 1} \smallsetminus  \{ 0\}  \rightarrow  {\mathbb{{RP}}}^{n}$ sending each point $x \in  {\mathbb{R}}^{n + 1} \smallsetminus  \{ 0\}$ to the subspace spanned by $x$ . The 2-dimensional projective space ${\mathbb{{RP}}}^{2}$ is called the projective plane. For any point $x \in  {\mathbb{R}}^{n + 1} \smallsetminus  \{ 0\}$ , let $\left\lbrack  x\right\rbrack   = \pi \left( x\right)  \in  {\mathbb{{RP}}}^{n}$ denote the line spanned by $x$ .

For each $i = 1,\ldots , n + 1$ , let ${\widetilde{U}}_{i} \subseteq  {\mathbb{R}}^{n + 1} \smallsetminus  \{ 0\}$ be the set where ${x}^{i} \neq  0$ , and let ${U}_{i} = \pi \left( {\widetilde{U}}_{i}\right)  \subseteq  {\mathbb{{RP}}}^{n}$ . Since ${\widetilde{U}}_{i}$ is a saturated open subset, ${U}_{i}$ is open and ${\left. \pi \right| }_{{\widetilde{U}}_{i}} : {\widetilde{U}}_{i} \rightarrow  {U}_{i}$ is a quotient map (see Theorem A.27). Define a map ${\varphi }_{i} : {U}_{i} \rightarrow  {\mathbb{R}}^{n}$ by

\[
{\varphi }_{i}\left\lbrack  {{x}^{1},\ldots ,{x}^{n + 1}}\right\rbrack   = \left( {\frac{{x}^{1}}{{x}^{i}},\ldots ,\frac{{x}^{i - 1}}{{x}^{i}},\frac{{x}^{i + 1}}{{x}^{i}},\ldots ,\frac{{x}^{n + 1}}{{x}^{i}}}\right) .
\]

This map is well defined because its value is unchanged by multiplying $x$ by a nonzero constant. Because ${\varphi }_{i} \circ  \pi$ is continuous, ${\varphi }_{i}$ is continuous by the characteristic property of quotient maps (Theorem A.27). In fact, ${\varphi }_{i}$ is a homeomorphism, because it has a continuous inverse given by

\[
{\varphi }_{i}^{-1}\left( {{u}^{1},\ldots ,{u}^{n}}\right)  = \left\lbrack  {{u}^{1},\ldots ,{u}^{i - 1},1,{u}^{i},\ldots ,{u}^{n}}\right\rbrack  ,
\]

as you can check. Geometrically, $\varphi \left( \left\lbrack  x\right\rbrack  \right)  = u$ means $\left( {u,1}\right)$ is the point in ${\mathbb{R}}^{n + 1}$ where the line $\left\lbrack  x\right\rbrack$ intersects the affine hyperplane where ${x}^{i} = 1$ (Fig. 1.4). Because the sets ${U}_{1},\ldots ,{U}_{n + 1}$ cover ${\mathbb{{RP}}}^{n}$ , this shows that ${\mathbb{{RP}}}^{n}$ is locally Euclidean of dimension $n$ . The Hausdorff and second-countability properties are left as exercises.

![21_374_191_784_436_0.jpg](images/21_374_191_784_436_0.jpg)

Fig. 1.4 A chart for ${\mathbb{{RP}}}^{n}$
\end{example}

\begin{exercise}
\label{exer:1-6}
\lean{realProjectiveSpace_chartDomain_isOpen}
\leanok
\dcref{ch1:1.6}
Show that ${\mathbb{{RP}}}^{n}$ is Hausdorff and second-countable, and is therefore a topological $n$ -manifold.
\end{exercise}

\begin{exercise}
\label{exer:1-7}
\lean{realProjectiveSpaceCompactSpace}
\leanok
\dcref{ch1:1.7}
Show that ${\mathbb{{RP}}}^{n}$ is compact. [Hint: show that the restriction of $\pi$ to ${\mathbb{S}}^{n}$ is surjective.]
\end{exercise}

\begin{example}[Product Manifolds]
\label{ex:1-8}
\lean{isManifold_pi}
\leanok
\dcref{ch1:1.8}
Suppose ${M}_{1},\ldots ,{M}_{k}$ are topological manifolds of dimensions ${n}_{1},\ldots ,{n}_{k}$ , respectively. The product space ${M}_{1} \times  \cdots  \times  {M}_{k}$ is shown to be a topological manifold of dimension ${n}_{1} + \cdots  + {n}_{k}$ as follows. It is Hausdorff and second-countable by Propositions A. 17 and A.23, so only the locally Euclidean property needs to be checked. Given any point $\left( {{p}_{1},\ldots ,{p}_{k}}\right)  \in \; {M}_{1} \times  \cdots  \times  {M}_{k}$ , we can choose a coordinate chart $\left( {{U}_{i},{\varphi }_{i}}\right)$ for each ${M}_{i}$ with ${p}_{i} \in  {U}_{i}$ . The product map

\[
{\varphi }_{1} \times  \cdots  \times  {\varphi }_{k} : {U}_{1} \times  \cdots  \times  {U}_{k} \rightarrow  {\mathbb{R}}^{{n}_{1} + \cdots  + {n}_{k}}
\]

is a homeomorphism onto its image, which is a product open subset of ${\mathbb{R}}^{{n}_{1} + \cdots  + {n}_{k}}$ . Thus, ${M}_{1} \times  \cdots  \times  {M}_{k}$ is a topological manifold of dimension ${n}_{1} + \cdots  + {n}_{k}$ , with charts of the form $\left( {{U}_{1} \times  \cdots  \times  {U}_{k},{\varphi }_{1} \times  \cdots  \times  {\varphi }_{k}}\right)$ .
\end{example}

\begin{example}[Tori]
\label{ex:1-9}
\lean{n_torus}
\leanok
\dcref{ch1:1.9}
For a positive integer $n$ , the $n$ -torus (plural: tori) is the product space ${\mathbb{T}}^{n} = {\mathbb{S}}^{1} \times  \cdots  \times  {\mathbb{S}}^{1}$ . By the discussion above, it is a topological $n$ -manifold. (The 2-torus is usually called simply the torus.)
\end{example}

\begin{definition}
\label{def:1-extra-3}
\lean{LocPathConnectedSpace}
\mathlibok
Recall that a topological space $X$ is

\begin{itemize}
\item connected if there do not exist two disjoint, nonempty, open subsets of $X$ whose union is $X$ ;
\end{itemize}

\begin{itemize}
\item path-connected if every pair of points in $X$ can be joined by a path in $X$ ; and
\end{itemize}

\begin{itemize}
\item locally path-connected if $X$ has a basis of path-connected open subsets.
\end{itemize}

(See Appendix A.)
\end{definition}

\begin{definition}
\label{def:1-extra-4}
\lean{TopologicalSpace.IsOpenCover}
\mathlibok
Let $M$ be a topological space. A collection $\mathcal{X}$ of subsets of $M$ is said to be locally finite if each point of $M$ has a neighborhood that intersects at most finitely many of the sets in $\mathcal{X}$ . Given a cover $\mathcal{U}$ of $M$ , another cover $\mathcal{V}$ is called a refinement of $\mathcal{U}$ if for each $V \in  \mathcal{V}$ there exists some $U \in  \mathcal{U}$ such that $V \subseteq  U$ . We say that $M$ is paracompact if every open cover of $M$ admits an open, locally finite refinement.
\end{definition}

\begin{lemma}
\label{lem:1-13}
\lean{LocallyFinite.closure_iUnion}
\mathlibok
\dcref{ch1:1.13}
Suppose $X$ is a locally finite collection of subsets of a topological space $M$ .

(a) The collection $\{ \bar{X} : X \in  \mathcal{X}\}$ is also locally finite.

(b) $\overline{\mathop{\bigcup }\limits_{{X \in  \mathcal{X}}}X} = \mathop{\bigcup }\limits_{{X \in  \mathcal{X}}}\bar{X}$ .
\end{lemma}

\begin{exercise}
\label{exer:1-14}
\lean{LocallyFinite.closure_iUnion}
\mathlibok
\uses{lem:1-13}
\dcref{ch1:1.14}
\begin{itemize}
\item Exercise 1.14. Prove the preceding lemma.
\end{itemize}
\end{exercise}


\section{Smooth Structures}

\begin{definition}
\label{def:1-2-extra-1}
\lean{Diffeomorph.toHomeomorph}
\mathlibok
If $U$ and $V$ are open subsets of Euclidean spaces ${\mathbb{R}}^{n}$ and ${\mathbb{R}}^{m}$ , respectively, a function $F : U \rightarrow  V$ is said to be smooth (or ${\mathbf{C}}^{\infty }$ , or infinitely differentiable) if each of its component functions has continuous partial derivatives of all orders. If in addition $F$ is bijective and has a smooth inverse map, it is called a diffeomorphism. A diffeomorphism is, in particular, a homeomorphism.
\end{definition}

\begin{definition}
\label{def:1-2-extra-2}
\lean{e.symm}
\mathlibok
Let $M$ be a topological $n$ -manifold. If $\left( {U,\varphi }\right) ,\left( {V,\psi }\right)$ are two charts such that $U \cap  V \neq  \varnothing$ , the composite map $\psi  \circ  {\varphi }^{-1} : \varphi \left( {U \cap  V}\right)  \rightarrow  \psi \left( {U \cap  V}\right)$ is called the transition map from $\varphi$ to $\psi$ (Fig. 1.6). It is a composition of homeomorphisms, and is therefore itself a homeomorphism. Two charts $\left( {U,\varphi }\right)$ and $\left( {V,\psi }\right)$ are said to be smoothly compatible if either $U \cap  V = \varnothing$ or the transition map $\psi  \circ  {\varphi }^{-1}$ is a diffeomorphism. Since $\varphi \left( {U \cap  V}\right)$ and $\psi \left( {U \cap  V}\right)$ are open subsets of ${\mathbb{R}}^{n}$ , smoothness of this map is to be interpreted in the ordinary sense of having continuous partial derivatives of all orders.
\end{definition}

\begin{definition}
\label{def:1-2-extra-3}
\lean{OpenPartialHomeomorph.IsSmoothAtlas}
\leanok
We define an atlas for $\mathbf{M}$ to be a collection of charts whose domains cover $M$ . An atlas $\mathcal{A}$ is called a smooth atlas if any two charts in $\mathcal{A}$ are smoothly compatible with each other.
\end{definition}

\begin{definition}
\label{def:1-2-extra-4}
\lean{StructureGroupoid.maximalAtlas}
\mathlibok
a smooth atlas $\mathcal{A}$ on $M$ is maximal if it is not properly contained in any larger smooth atlas. This just means that any chart that is smoothly compatible with every chart in $\mathcal{A}$ is already in $\mathcal{A}$ . (Such a smooth atlas is also said to be complete.)
\end{definition}

\begin{definition}
\label{def:1-2-extra-5}
\lean{smooth_manifold_smooth_structure}
\leanok
If $M$ is a topological manifold, a smooth structure on $\mathbf{M}$ is a maximal smooth atlas. A smooth manifold is a pair $\left( {M,\mathcal{A}}\right)$ , where $M$ is a topological manifold and $\mathcal{A}$ is a smooth structure on $M$ . When the smooth structure is understood, we usually omit mention of it and just say " $M$ is a smooth manifold." Smooth structures are also called differentiable structures or ${C}^{\infty }$ structures by some authors. We also use the term smooth manifold structure to mean a manifold topology together with a smooth structure.
\end{definition}

\begin{proposition}
\label{prop:1-17}
\lean{same_smooth_structure_iff_union_is_smooth_atlas}
\leanok
\dcref{ch1:1.17}
Let $M$ be a topological manifold.

(a) Every smooth atlas $\mathcal{A}$ for $M$ is contained in a unique maximal smooth atlas, called the smooth structure determined by $\mathcal{A}$ .

(b) Two smooth atlases for $M$ determine the same smooth structure if and only if their union is a smooth atlas.
\end{proposition}

\begin{proof}
\leanok
Let $\mathcal{A}$ be a smooth atlas for $M$ , and let $\overline{\mathcal{A}}$ denote the set of all charts that are smoothly compatible with every chart in $\mathcal{A}$ . To show that $\overline{\mathcal{A}}$ is a smooth atlas, we need to show that any two charts of $\overline{\mathcal{A}}$ are smoothly compatible with each other, which is to say that for any $\left( {U,\varphi }\right) ,\left( {V,\psi }\right)  \in  \overline{\mathcal{A}}$ , the map $\psi  \circ  {\varphi }^{-1} : \varphi \left( {U \cap  V}\right)  \rightarrow \; \psi \left( {U \cap  V}\right)$ is smooth.

Let $x = \varphi \left( p\right)  \in  \varphi \left( {U \cap  V}\right)$ be arbitrary. Because the domains of the charts in $\mathcal{A}$ cover $M$ , there is some chart $\left( {W,\theta }\right)  \in  \mathcal{A}$ such that $p \in  W$ (Fig. 1.7). Since every chart in $\overline{\mathcal{A}}$ is smoothly compatible with $\left( {W,\theta }\right)$ , both of the maps $\theta  \circ  {\varphi }^{-1}$ and $\psi  \circ  {\theta }^{-1}$ are smooth where they are defined. Since $p \in  U \cap  V \cap  W$ , it follows that $\psi  \circ  {\varphi }^{-1} = \; \left( {\psi  \circ  {\theta }^{-1}}\right)  \circ  \left( {\theta  \circ  {\varphi }^{-1}}\right)$ is smooth on a neighborhood of $x$ . Thus, $\psi  \circ  {\varphi }^{-1}$ is smooth in a neighborhood of each point in $\varphi \left( {U \cap  V}\right)$ . Therefore, $\overline{\mathcal{A}}$ is a smooth atlas. To check that it is maximal, just note that any chart that is smoothly compatible with every chart in $\overline{\mathcal{A}}$ must in particular be smoothly compatible with every chart in $\mathcal{A}$ , so it is already in $\overline{\mathcal{A}}$ . This proves the existence of a maximal smooth atlas containing $\mathcal{A}$ . If $\mathcal{B}$ is any other maximal smooth atlas containing $\mathcal{A}$ , each of its charts is smoothly compatible with each chart in $\mathcal{A}$ , so $\mathcal{B} \subseteq  \overline{\mathcal{A}}$ . By maximality of $\mathcal{B},\mathcal{B} = \overline{\mathcal{A}}$ .

The proof of (b) is left as an exercise.
\end{proof}

\begin{exercise}
\label{exer:1-18}
\lean{same_smooth_structure_iff_union_is_smooth_atlas}
\leanok
\uses{prop:1-17}
\dcref{ch1:1.18}
\begin{itemize}
\item Exercise 1.18. Prove Proposition 1.17(b).
\end{itemize}
\end{exercise}

\begin{remark}
\label{rem:1-2-extra-6}
\lean{IsManifold}
\mathlibok
It is worth mentioning that the notion of smooth structure can be generalized in several different ways by changing the compatibility requirement for charts. For example, if we replace the requirement that charts be smoothly compatible by the weaker requirement that each transition map $\psi  \circ  {\varphi }^{-1}$ (and its inverse) be of class ${C}^{k}$ , we obtain the definition of a ${\mathbf{C}}^{k}$ structure. Similarly, if we require that each transition map be real-analytic (i.e., expressible as a convergent power series in a neighborhood of each point), we obtain the definition of a real-analytic structure, also called a ${\mathbf{C}}^{\omega }$ structure. If $M$ has even dimension $n = {2m}$ , we can identify ${\mathbb{R}}^{2m}$ with ${\mathbb{C}}^{m}$ and require that the transition maps be complex-analytic; this determines a complex-analytic structure. A manifold endowed with one of these structures is called a ${C}^{k}$ manifold, real-analytic manifold, or complex manifold, respectively. (Note that a ${C}^{0}$ manifold is just a topological manifold.) We do not treat any of these other kinds of manifolds in this book, but they play important roles in analysis, so it is useful to know the definitions.
\end{remark}


\section{Local Coordinate Representations}

\begin{definition}
\label{def:1-3-extra-1}
\lean{IsRegularCoordinateBall}
\leanok
If $M$ is a smooth manifold, any chart $\left( {U,\varphi }\right)$ contained in the given maximal smooth atlas is called a smooth chart, and the corresponding coordinate map $\varphi$ is called a smooth coordinate map. It is useful also to introduce the terms smooth coordinate domain or smooth coordinate neighborhood for the domain of a smooth coordinate chart. A smooth coordinate ball means a smooth coordinate domain whose image under a smooth coordinate map is a ball in Euclidean space. A smooth coordinate cube is defined similarly.

It is often useful to restrict attention to coordinate balls whose closures sit nicely inside larger coordinate balls. We say a set $B \subseteq  M$ is a regular coordinate ball if there is a smooth coordinate ball ${B}^{\prime } \supseteq  \bar{B}$ and a smooth coordinate map $\varphi  : {B}^{\prime } \rightarrow  {\mathbb{R}}^{n}$ such that for some positive real numbers $r < {r}^{\prime }$ ,

\[
\varphi \left( B\right)  = {B}_{r}\left( 0\right) ,\;\varphi \left( \bar{B}\right)  = {\bar{B}}_{r}\left( 0\right) ,\;\text{ and }\;\varphi \left( {B}^{\prime }\right)  = {B}_{{r}^{\prime }}\left( 0\right) .
\]
\end{definition}

\begin{proposition}
\label{prop:1-19}
\lean{exists_countable_regular_coordinate_ball_basis}
\leanok
\dcref{ch1:1.19}
Every smooth manifold has a countable basis of regular coordinate balls.
\end{proposition}

\begin{exercise}
\label{exer:1-20}
\lean{exists_countable_regular_coordinate_ball_basis}
\leanok
\dcref{ch1:1.20}
Prove Proposition 1.19.
\end{exercise}

\begin{notation}
\label{not:1-3-extra-2}
\lean{extChartAt}
\mathlibok
Here is how one usually thinks about coordinate charts on a smooth manifold. Once we choose a smooth chart $\left( {U,\varphi }\right)$ on $M$ , the coordinate map $\varphi  : U \rightarrow  \widehat{U} \subseteq  {\mathbb{R}}^{n}$ can be thought of as giving a temporary identification between $U$ and $\widehat{U}$ . Using this identification, while we work in this chart, we can think of $U$ simultaneously as an open subset of $M$ and as an open subset of ${\mathbb{R}}^{n}$ . You can visualize this identification by thinking of a "grid" drawn on $U$ representing the preimages of the coordinate lines under $\varphi$ (Fig. 1.8). Under this identification, we can represent a point $p \in \; U$ by its coordinates $\left( {{x}^{1},\ldots ,{x}^{n}}\right)  = \varphi \left( p\right)$ , and think of this $n$ -tuple as being the point $p$ . We typically express this by saying " $\left( {{x}^{1},\ldots ,{x}^{n}}\right)$ is the (local) coordinate representation for $p$ " or " $p = \left( {{x}^{1},\ldots ,{x}^{n}}\right)$ in local coordinates."
\end{notation}


\section{Examples of Smooth Manifolds}

\begin{example}[Euclidean Spaces]
\label{ex:1-22}
\lean{euclideanSpace_standard_smooth_structure}
\leanok
\dcref{ch1:1.22}
For each nonnegative integer $n$ , the Euclidean space ${\mathbb{R}}^{n}$ is a smooth $n$ -manifold with the smooth structure determined by the atlas consisting of the single chart $\left( {{\mathbb{R}}^{n},{\operatorname{Id}}_{{\mathbb{R}}^{n}}}\right)$ . We call this the standard smooth structure on ${\mathbb{R}}^{n}$ and the resulting coordinate map standard coordinates. Unless we explicitly specify otherwise, we always use this smooth structure on ${\mathbb{R}}^{n}$ . With respect to this smooth structure, the smooth coordinate charts for ${\mathbb{R}}^{n}$ are exactly those charts $\left( {U,\varphi }\right)$ such that $\varphi$ is a diffeomorphism (in the sense of ordinary calculus) from $U$ to another open subset $\widehat{U} \subseteq  {\mathbb{R}}^{n}$ .
\end{example}

\begin{example}[Another Smooth Structure on $\mathbb{R}$]
\label{ex:1-23}
\lean{cubicChart_not_mem_standard_smooth_maximalAtlas}
\leanok
\dcref{ch1:1.23}
Consider the homeomorphism $\psi  : \mathbb{R} \rightarrow  \mathbb{R}$ given by

\[
\psi \left( x\right)  = {x}^{3}. \tag{1.1}
\]

The atlas consisting of the single chart $\left( {\mathbb{R},\psi }\right)$ defines a smooth structure on $\mathbb{R}$ . This chart is not smoothly compatible with the standard smooth structure, because the transition map ${\operatorname{Id}}_{\mathbb{R}} \circ  {\psi }^{-1}\left( y\right)  = {y}^{1/3}$ is not smooth at the origin. Therefore, the smooth structure defined on $\mathbb{R}$ by $\psi$ is not the same as the standard one. Using similar ideas, it is not hard to construct many distinct smooth structures on any given positive-dimensional topological manifold, as long as it has one smooth structure to begin with (see Problem 1-6).
\end{example}

\begin{example}[Finite-Dimensional Vector Spaces]
\label{ex:1-24}
\lean{basis_coordinate_change_contDiff}
\leanok
\dcref{ch1:1.24}
Let $V$ be a finite-dimensional real vector space. Any norm on $V$ determines a topology, which is independent of the choice of norm (Exercise B.49). With this topology, $V$ is a topological $n$ - manifold, and has a natural smooth structure defined as follows. Each (ordered) basis $\left( {{E}_{1},\ldots ,{E}_{n}}\right)$ for $V$ defines a basis isomorphism $E : {\mathbb{R}}^{n} \rightarrow  V$ by

\[
E\left( x\right)  = \mathop{\sum }\limits_{{i = 1}}^{n}{x}^{i}{E}_{i}
\]

This map is a homeomorphism, so $\left( {V,{E}^{-1}}\right)$ is a chart. If $\left( {{\widetilde{E}}_{1},\ldots ,{\widetilde{E}}_{n}}\right)$ is any other basis and $\widetilde{E}\left( x\right)  = \mathop{\sum }\limits_{j}{x}^{j}{\widetilde{E}}_{j}$ is the corresponding isomorphism, then there is some invertible matrix $\left( {A}_{i}^{j}\right)$ such that ${E}_{i} = \mathop{\sum }\limits_{j}{A}_{i}^{j}{\widetilde{E}}_{j}$ for each $i$ . The transition map between the two charts is then given by ${\widetilde{E}}^{-1} \circ  E\left( x\right)  = \widetilde{x}$ , where $\widetilde{x} = \left( {{\widetilde{x}}^{1},\ldots ,{\widetilde{x}}^{n}}\right)$ is determined by

\[
\mathop{\sum }\limits_{{j = 1}}^{n}{\widetilde{x}}^{j}{\widetilde{E}}_{j} = \mathop{\sum }\limits_{{i = 1}}^{n}{x}^{i}{E}_{i} = \mathop{\sum }\limits_{{i, j = 1}}^{n}{x}^{i}{A}_{i}^{j}{\widetilde{E}}_{j}.
\]

It follows that ${\widetilde{x}}^{j} = \mathop{\sum }\limits_{i}{A}_{i}^{j}{x}^{i}$ . Thus, the map sending $x$ to $\widetilde{x}$ is an invertible linear map and hence a diffeomorphism, so any two such charts are smoothly compatible. The collection of all such charts thus defines a smooth structure, called the standard smooth structure on $V$ .
\end{example}

\begin{notation}
\label{not:1-4-extra-1}
\lean{Basis.sum_repr}
\mathlibok
This is a good place to pause and introduce an important notational convention that is commonly used in the study of smooth manifolds. Because of the proliferation of summations such as $\mathop{\sum }\limits_{i}{x}^{i}{E}_{i}$ in this subject, we often abbreviate such a sum by omitting the summation sign, as in

\[
E\left( x\right)  = {x}^{i}{E}_{i},\;\text{ an abbreviation for }E\left( x\right)  = \mathop{\sum }\limits_{{i = 1}}^{n}{x}^{i}{E}_{i}.
\]

We interpret any such expression according to the following rule, called the Einstein summation convention: if the same index name (such as $i$ in the expression above) appears exactly twice in any monomial term, once as an upper index and once as a lower index, that term is understood to be summed over all possible values of that index, generally from 1 to the dimension of the space in question. This simple idea was introduced by Einstein to reduce the complexity of expressions arising in the study of smooth manifolds by eliminating the necessity of explicitly writing summation signs. We use the summation convention systematically throughout the book (except in the appendices, which many readers will look at before the rest of the book).

Another important aspect of the summation convention is the positions of the indices. We always write basis vectors (such as ${E}_{i}$ ) with lower indices, and components of a vector with respect to a basis (such as ${x}^{i}$ ) with upper indices. These index conventions help to ensure that, in summations that make mathematical sense, each index to be summed over typically appears twice in any given term, once as a lower index and once as an upper index. Any index that is implicitly summed over is a "dummy index," meaning that the value of such an expression is unchanged if a different name is substituted for each dummy index. For example, ${x}^{i}{E}_{i}$ and ${x}^{j}{E}_{j}$ mean exactly the same thing.

Since the coordinates of a point $\left( {{x}^{1},\ldots ,{x}^{n}}\right)  \in  {\mathbb{R}}^{n}$ are also its components with respect to the standard basis, in order to be consistent with our convention of writing components of vectors with upper indices, we need to use upper indices for these coordinates, and we do so throughout this book. Although this may seem awkward at first, in combination with the summation convention it offers enormous advantages when we work with complicated indexed sums, not the least of which is that expressions that are not mathematically meaningful often betray themselves quickly by violating the index convention. (The main exceptions are expressions involving the Euclidean dot product $x \cdot  y = \mathop{\sum }\limits_{i}{x}^{i}{y}^{i}$ , in which the same index appears twice in the upper position, and the standard symplectic form on ${\mathbb{R}}^{2n}$ , which we will define in Chapter 22. We always explicitly write summation signs in such expressions.)
\end{notation}

\begin{example}[The General Linear Group]
\label{ex:1-27}
\lean{Units.isOpenEmbedding_val}
\mathlibok
\dcref{ch1:1.27}
The general linear group $\mathrm{{GL}}\left( {n,\mathbb{R}}\right)$ is the set of invertible $n \times  n$ matrices with real entries. It is a smooth ${n}^{2}$ -dimensional manifold because it is an open subset of the ${n}^{2}$ -dimensional vector space $\mathrm{M}\left( {n,\mathbb{R}}\right)$ , namely the set where the (continuous) determinant function is nonzero.
\end{example}

\begin{example}[Matrices of Full Rank]
\label{ex:1-28}
\lean{fullRankRange_iff_natLeRank}
\leanok
\dcref{ch1:1.28}
The previous example has a natural generalization to rectangular matrices of full rank. Suppose $m < n$ , and let ${\mathrm{M}}_{m}\left( {m \times  n,\mathbb{R}}\right)$ denote the subset of $\mathrm{M}\left( {m \times  n,\mathbb{R}}\right)$ consisting of matrices of rank $m$ . If $A$ is an arbitrary such matrix, the fact that $\operatorname{rank}A = m$ means that $A$ has some nonsingular $m \times  m$ submatrix. By continuity of the determinant function, this same submatrix has nonzero determinant on a neighborhood of $A$ in $\mathrm{M}\left( {m \times  n,\mathbb{R}}\right)$ , which implies that $A$ has a neighborhood contained in ${\mathrm{M}}_{m}\left( {m \times  n,\mathbb{R}}\right)$ . Thus, ${\mathrm{M}}_{m}\left( {m \times  n,\mathbb{R}}\right)$ is an open subset of $\mathrm{M}\left( {m \times  n,\mathbb{R}}\right)$ , and therefore is itself a smooth ${mn}$ -dimensional manifold. A similar argument shows that ${\mathrm{M}}_{n}\left( {m \times  n,\mathbb{R}}\right)$ is a smooth ${mn}$ -manifold when $n < m$ .
\end{example}

\begin{example}[Graphs of Smooth Functions]
\label{ex:1-30}
\lean{graph_smooth_structure}
\leanok
\uses{ex:1-3}
\dcref{ch1:1.30}
If $U \subseteq  {\mathbb{R}}^{n}$ is an open subset and $f : U \rightarrow  {\mathbb{R}}^{k}$ is a smooth function, we have already observed above (Example 1.3) that the graph of $f$ is a topological $n$ -manifold in the subspace topology. Since $\Gamma \left( f\right)$ is covered by the single graph coordinate chart $\varphi  : \Gamma \left( f\right)  \rightarrow  U$ (the restriction of ${\pi }_{1}$ ), we can put a canonical smooth structure on $\Gamma \left( f\right)$ by declaring the graph coordinate chart $\left( {\Gamma \left( f\right) ,\varphi }\right)$ to be a smooth chart.
\end{example}

\begin{example}[Projective Spaces]
\label{ex:1-33}
\lean{realProjectiveChartTransitionDiffeomorph}
\leanok
\uses{ex:1-5}
\dcref{ch1:1.33}
The $n$ -dimensional real projective space ${\mathbb{{RP}}}^{n}$ is a topological $n$ -manifold by Example 1.5. Let us check that the coordinate charts $\left( {{U}_{i},{\varphi }_{i}}\right)$ constructed in that example are all smoothly compatible. Assuming for convenience that $i > j$ , it is straightforward to compute that

\[
{\varphi }_{j} \circ  {\varphi }_{i}^{-1}\left( {{u}^{1},\ldots ,{u}^{n}}\right)  = \left( {\frac{{u}^{1}}{{u}^{j}},\ldots ,\frac{{u}^{j - 1}}{{u}^{j}},\frac{{u}^{j + 1}}{{u}^{j}},\ldots ,\frac{{u}^{i - 1}}{{u}^{j}},\frac{1}{{u}^{j}},\frac{{u}^{i}}{{u}^{j}},\ldots ,\frac{{u}^{n}}{{u}^{j}}}\right) ,
\]

which is a diffeomorphism from ${\varphi }_{i}\left( {{U}_{i} \cap  {U}_{j}}\right)$ to ${\varphi }_{j}\left( {{U}_{i} \cap  {U}_{j}}\right)$ .
\end{example}

\begin{lemma}[Smooth Manifold Chart Lemma]
\label{lem:1-35}
\lean{ChartedSpaceCore.eq_toTopologicalSpace_of_chartedSpace}
\leanok
\dcref{ch1:1.35}
Let $M$ be a set, and suppose we are given a collection $\left\{  {U}_{\alpha }\right\}$ of subsets of $M$ together with maps ${\varphi }_{\alpha } : {U}_{\alpha } \rightarrow  {\mathbb{R}}^{n}$ , such that the following properties are satisfied:

(i) For each $\alpha ,{\varphi }_{\alpha }$ is a bijection between ${U}_{\alpha }$ and an open subset ${\varphi }_{\alpha }\left( {U}_{\alpha }\right)  \subseteq  {\mathbb{R}}^{n}$ .

(ii) For each $\alpha$ and $\beta$ , the sets ${\varphi }_{\alpha }\left( {{U}_{\alpha } \cap  {U}_{\beta }}\right)$ and ${\varphi }_{\beta }\left( {{U}_{\alpha } \cap  {U}_{\beta }}\right)$ are open in ${\mathbb{R}}^{n}$ .

(iii) Whenever ${U}_{\alpha } \cap  {U}_{\beta } \neq  \varnothing$ , the map ${\varphi }_{\beta } \circ  {\varphi }_{\alpha }^{-1} : {\varphi }_{\alpha }\left( {{U}_{\alpha } \cap  {U}_{\beta }}\right)  \rightarrow  {\varphi }_{\beta }\left( {{U}_{\alpha } \cap  {U}_{\beta }}\right)$ is smooth.

(iv) Countably many of the sets ${U}_{\alpha }$ cover $M$ .

(v) Whenever $p, q$ are distinct points in $M$ , either there exists some ${U}_{\alpha }$ containing both $p$ and $q$ or there exist disjoint sets ${U}_{\alpha },{U}_{\beta }$ with $p \in  {U}_{\alpha }$ and $q \in  {U}_{\beta }$ .

Then $M$ has a unique smooth manifold structure such that each $\left( {{U}_{\alpha },{\varphi }_{\alpha }}\right)$ is a smooth chart.

![36_269_200_1002_699_0.jpg](images/36_269_200_1002_699_0.jpg)

Fig. 1.9 The smooth manifold chart lemma
\end{lemma}

\begin{proof}
\leanok
We define the topology by taking all sets of the form ${\varphi }_{\alpha }^{-1}\left( V\right)$ , with $V$ an open subset of ${\mathbb{R}}^{n}$ , as a basis. To prove that this is a basis for a topology, we need to show that for any point $p$ in the intersection of two basis sets ${\varphi }_{\alpha }^{-1}\left( V\right)$ and ${\varphi }_{\beta }^{-1}\left( W\right)$ , there is a third basis set containing $p$ and contained in the intersection. It suffices to show that ${\varphi }_{\alpha }^{-1}\left( V\right)  \cap  {\varphi }_{\beta }^{-1}\left( W\right)$ is itself a basis set (Fig. 1.9). To see this, observe that (iii) implies that ${\left( {\varphi }_{\beta } \circ  {\varphi }_{\alpha }^{-1}\right) }^{-1}\left( W\right)$ is an open subset of ${\varphi }_{\alpha }\left( {{U}_{\alpha } \cap  {U}_{\beta }}\right)$ , and (ii) implies that this set is also open in ${\mathbb{R}}^{n}$ . It follows that

\[
{\varphi }_{\alpha }^{-1}\left( V\right)  \cap  {\varphi }_{\beta }^{-1}\left( W\right)  = {\varphi }_{\alpha }^{-1}\left( {V \cap  {\left( {\varphi }_{\beta } \circ  {\varphi }_{\alpha }^{-1}\right) }^{-1}\left( W\right) }\right)
\]

is also a basis set, as claimed.

Each map ${\varphi }_{\alpha }$ is then a homeomorphism onto its image (essentially by definition), so $M$ is locally Euclidean of dimension $n$ . The Hausdorff property follows easily from (v), and second-countability follows from (iv) and the result of Exercise A.22, because each ${U}_{\alpha }$ is second-countable. Finally,(iii) guarantees that the collection $\left\{  \left( {{U}_{\alpha },{\varphi }_{\alpha }}\right) \right\}$ is a smooth atlas. It is clear that this topology and smooth structure are the unique ones satisfying the conclusions of the lemma.
\end{proof}

\begin{example}[Grassmann Manifolds]
\label{ex:1-36}
\lean{grassmannian.chart_equiv}
\leanok
\uses{ex:1-28,lem:1-35}
\dcref{ch1:1.36}
Let $V$ be an $n$ -dimensional real vector space. For any integer $0 \leq  k \leq  n$ , we let ${\mathrm{G}}_{k}\left( V\right)$ denote the set of all $k$ -dimensional linear subspaces of $V$ . We will show that ${\mathrm{G}}_{k}\left( V\right)$ can be naturally given the structure of a smooth manifold of dimension $k\left( {n - k}\right)$ . With this structure, it is called a Grassmann manifold, or simply a Grassmannian. In the special case $V = {\mathbb{R}}^{n}$ , the Grassmannian ${\mathrm{G}}_{k}\left( {\mathbb{R}}^{n}\right)$ is often denoted by some simpler notation such as ${\mathrm{G}}_{k, n}$ or $\mathrm{G}\left( {k, n}\right)$ . Note that ${\mathrm{G}}_{1}\left( {\mathbb{R}}^{n + 1}\right)$ is exactly the $n$ -dimensional projective space ${\mathbb{{RP}}}^{n}$ .

The construction of a smooth structure on ${\mathrm{G}}_{k}\left( V\right)$ is somewhat more involved than the ones we have done so far, but the basic idea is just to use linear algebra to construct charts for ${\mathrm{G}}_{k}\left( V\right)$ , and then apply the smooth manifold chart lemma. We will give a shorter proof that ${\mathrm{G}}_{k}\left( V\right)$ is a smooth manifold in Chapter 21 (see Example 21.21).

Let $P$ and $Q$ be any complementary subspaces of $V$ of dimensions $k$ and $n - k$ , respectively, so that $V$ decomposes as a direct sum: $V = P \oplus  Q$ . The graph of any linear map $X : P \rightarrow  Q$ can be identified with a $k$ -dimensional subspace $\Gamma \left( X\right)  \subseteq  V$ , defined by

\[
\Gamma \left( X\right)  = \{ v + {Xv} : v \in  P\} .
\]

Any such subspace has the property that its intersection with $Q$ is the zero subspace. Conversely, any subspace $S \subseteq  V$ that intersects $Q$ trivially is the graph of a unique linear map $X : P \rightarrow  Q$ , which can be constructed as follows: let ${\pi }_{P} : V \rightarrow  P$ and ${\pi }_{Q} : V \rightarrow  Q$ be the projections determined by the direct sum decomposition; then the hypothesis implies that ${\left. {\pi }_{P}\right| }_{S}$ is an isomorphism from $S$ to $P$ . Therefore, $X = \left( {\left. {\pi }_{Q}\right| }_{S}\right)  \circ  {\left( {\left. {\pi }_{P}\right| }_{S}\right) }^{-1}$ is a well-defined linear map from $P$ to $Q$ , and it is straightforward to check that $S$ is its graph.

Let $\mathrm{L}\left( {P;Q}\right)$ denote the vector space of linear maps from $P$ to $Q$ , and let ${U}_{Q}$ denote the subset of ${\mathrm{G}}_{k}\left( V\right)$ consisting of $k$ -dimensional subspaces whose intersections with $Q$ are trivial. The assignment $X \mapsto  \Gamma \left( X\right)$ defines a map $\Gamma  : \mathrm{L}\left( {P;Q}\right)  \rightarrow  {U}_{Q}$ , and the discussion above shows that $\Gamma$ is a bijection. Let $\varphi  = {\Gamma }^{-1} : {U}_{Q} \rightarrow  \mathrm{L}\left( {P;Q}\right)$ . By choosing bases for $P$ and $Q$ , we can identify $\mathrm{L}\left( {P;Q}\right)$ with $\mathrm{M}\left( {\left( {n - k}\right)  \times  k,\mathbb{R}}\right)$ and hence with ${\mathbb{R}}^{k\left( {n - k}\right) }$ , and thus we can think of $\left( {{U}_{Q},\varphi }\right)$ as a coordinate chart. Since the image of each such chart is all of $\mathrm{L}\left( {P;Q}\right)$ , condition (i) of Lemma 1.35 is clearly satisfied.

Now let $\left( {{P}^{\prime },{Q}^{\prime }}\right)$ be any other such pair of subspaces, and let ${\pi }_{{P}^{\prime }},{\pi }_{{Q}^{\prime }}$ be the corresponding projections and ${\varphi }^{\prime } : {U}_{{Q}^{\prime }} \rightarrow  \mathrm{L}\left( {{P}^{\prime };{Q}^{\prime }}\right)$ the corresponding map. The set $\varphi \left( {{U}_{Q} \cap  {U}_{{Q}^{\prime }}}\right)  \subseteq  \mathrm{L}\left( {P;Q}\right)$ consists of all linear maps $X : P \rightarrow  Q$ whose graphs intersect ${Q}^{\prime }$ trivially. To see that this set is open in $\mathrm{L}\left( {P;Q}\right)$ , for each $X \in  \mathrm{L}\left( {P;Q}\right)$ , let ${I}_{X} : P \rightarrow  V$ be the map ${I}_{X}\left( v\right)  = v + {Xv}$ , which is a bijection from $P$ to the graph of $X$ . Because $\Gamma \left( X\right)  = \operatorname{Im}{I}_{X}$ and ${Q}^{\prime } = \operatorname{Ker}{\pi }_{{P}^{\prime }}$ , it follows from Exercise B. 22(d) that the graph of $X$ intersects ${Q}^{\prime }$ trivially if and only if ${\pi }_{{P}^{\prime }} \circ  {I}_{X}$ has full rank. Because the matrix entries of ${\pi }_{{P}^{\prime }} \circ  {I}_{X}$ (with respect to any bases) depend continuously on $X$ , the result of Example 1.28 shows that the set of all such $X$ is open in $\mathrm{L}\left( {P;Q}\right)$ . Thus property (ii) in the smooth manifold chart lemma holds.

We need to show that the transition map ${\varphi }^{\prime } \circ  {\varphi }^{-1}$ is smooth on $\varphi \left( {{U}_{Q} \cap  {U}_{{Q}^{\prime }}}\right)$ . Suppose $X \in  \varphi \left( {{U}_{Q} \cap  {U}_{{Q}^{\prime }}}\right)  \subseteq  \mathrm{L}\left( {P;Q}\right)$ is arbitrary, and let $S$ denote the subspace $\Gamma \left( X\right)  \subseteq  V$ . If we put ${X}^{\prime } = {\varphi }^{\prime } \circ  {\varphi }^{-1}\left( X\right)$ , then as above, ${X}^{\prime } = \left( {\left. {\pi }_{{Q}^{\prime }}\right| }_{S}\right)  \circ  {\left( {\left. {\pi }_{{P}^{\prime }}\right| }_{S}\right) }^{-1}$ (see Fig. 1.10). To relate this map to $X$ , note that ${I}_{X} : P \rightarrow  S$ is an isomorphism, so we can write

\[
{X}^{\prime } = \left( {\left. {\pi }_{{Q}^{\prime }}\right| }_{S}\right)  \circ  {I}_{X} \circ  {\left( {I}_{X}\right) }^{-1} \circ  {\left( {\left. {\pi }_{{P}^{\prime }}\right| }_{S}\right) }^{-1} = \left( {{\pi }_{{Q}^{\prime }} \circ  {I}_{X}}\right)  \circ  {\left( {\pi }_{{P}^{\prime }} \circ  {I}_{X}\right) }^{-1}.
\]

![38_524_187_487_474_0.jpg](images/38_524_187_487_474_0.jpg)

Fig. 1.10 Smooth compatibility of coordinates on ${\mathrm{G}}_{k}\left( V\right)$

To show that this depends smoothly on $X$ , define linear maps $A : P \rightarrow  {P}^{\prime }$ , $B : P \rightarrow  {Q}^{\prime }, C : Q \rightarrow  {P}^{\prime }$ , and $D : Q \rightarrow  {Q}^{\prime }$ as follows:

\[
A = {\left. {\pi }_{{P}^{\prime }}\right| }_{P},\;B = {\left. {\pi }_{{Q}^{\prime }}\right| }_{P},\;C = {\left. {\pi }_{{P}^{\prime }}\right| }_{Q},\;D = {\left. {\pi }_{{Q}^{\prime }}\right| }_{Q}.
\]

Then for $v \in  P$ , we have

\[
\left( {{\pi }_{{P}^{\prime }} \circ  {I}_{X}}\right) v = \left( {A + {CX}}\right) v,\;\left( {{\pi }_{{Q}^{\prime }} \circ  {I}_{X}}\right) v = \left( {B + {DX}}\right) v,
\]

from which it follows that ${X}^{\prime } = \left( {B + {DX}}\right) {\left( A + CX\right) }^{-1}$ . Once we choose bases for $P, Q,{P}^{\prime }$ , and ${Q}^{\prime }$ , all of these linear maps are represented by matrices. Because the matrix entries of ${\left( A + CX\right) }^{-1}$ are rational functions of those of $A + {CX}$ by Cramer’s rule, it follows that the matrix entries of ${X}^{\prime }$ depend smoothly on those of $X$ . This proves that ${\varphi }^{\prime } \circ  {\varphi }^{-1}$ is a smooth map, so the charts we have constructed satisfy condition (iii) of Lemma 1.35.

To check condition (iv), we just note that ${\mathrm{G}}_{k}\left( V\right)$ can in fact be covered by finitely many of the sets ${U}_{Q}$ : for example, if $\left( {{E}_{1},\ldots ,{E}_{n}}\right)$ is any fixed basis for $V$ , any partition of the basis elements into two subsets containing $k$ and $n - k$ elements determines appropriate subspaces $P$ and $Q$ , and any subspace $S$ must have trivial intersection with $Q$ for at least one of these partitions (see Exercise B.9). Thus, ${\mathrm{G}}_{k}\left( V\right)$ is covered by the finitely many charts determined by all possible partitions of a fixed basis.

Finally, the Hausdorff condition (v) is easily verified by noting that for any two $k$ - dimensional subspaces $P,{P}^{\prime } \subseteq  V$ , it is possible to find a subspace $Q$ of dimension $n - k$ whose intersections with both $P$ and ${P}^{\prime }$ are trivial, and then $P$ and ${P}^{\prime }$ are both contained in the domain of the chart determined by, say, $\left( {P, Q}\right)$ .
\end{example}


\section{Manifolds with Boundary}

\begin{definition}
\label{def:1-5-extra-1}
\lean{TopologicalManifoldWithBoundary}
\leanok
In many important applications of manifolds, most notably those involving integration, we will encounter spaces that would be smooth manifolds except that they have a "boundary" of some sort. Simple examples of such spaces include closed intervals in $\mathbb{R}$ , closed balls in ${\mathbb{R}}^{n}$ , and closed hemispheres in ${\mathbb{S}}^{n}$ . To accommodate such spaces, we need to extend our definition of manifolds.

![39_262_195_1006_423_0.jpg](images/39_262_195_1006_423_0.jpg)

Fig. 1.11 A manifold with boundary

Points in these spaces will have neighborhoods modeled either on open subsets of ${\mathbb{R}}^{n}$ or on open subsets of the closed n-dimensional upper half-space ${\mathbb{H}}^{n} \subseteq  {\mathbb{R}}^{n}$ , defined as

\[
{\mathbb{H}}^{n} = \left\{  {\left( {{x}^{1},\ldots ,{x}^{n}}\right)  \in  {\mathbb{R}}^{n} : {x}^{n} \geq  0}\right\}  .
\]

We will use the notations Int ${\mathbb{H}}^{n}$ and $\partial {\mathbb{H}}^{n}$ to denote the interior and boundary of ${\mathbb{H}}^{n}$ , respectively, as a subset of ${\mathbb{R}}^{n}$ . When $n > 0$ , this means

\[
\text{ Int }{\mathbb{H}}^{n} = \left\{  {\left( {{x}^{1},\ldots ,{x}^{n}}\right)  \in  {\mathbb{R}}^{n} : {x}^{n} > 0}\right\}  \text{ , }
\]

\[
\partial {\mathbb{H}}^{n} = \left\{  {\left( {{x}^{1},\ldots ,{x}^{n}}\right)  \in  {\mathbb{R}}^{n} : {x}^{n} = 0}\right\}  .
\]

In the $n = 0$ case, ${\mathbb{H}}^{0} = {\mathbb{R}}^{0} = \{ 0\}$ , so Int ${\mathbb{H}}^{0} = {\mathbb{R}}^{0}$ and $\partial {\mathbb{H}}^{0} = \varnothing$ .

An n-dimensional topological manifold with boundary is a second-countable Hausdorff space $M$ in which every point has a neighborhood homeomorphic either to an open subset of ${\mathbb{R}}^{n}$ or to a (relatively) open subset of ${\mathbb{H}}^{n}$ (Fig. 1.11). An open subset $U \subseteq  M$ together with a map $\varphi  : U \rightarrow  {\mathbb{R}}^{n}$ that is a homeomorphism onto an open subset of ${\mathbb{R}}^{n}$ or ${\mathbb{H}}^{n}$ will be called a chart for $M$ , just as in the case of manifolds. When it is necessary to make the distinction, we will call $\left( {U,\varphi }\right)$ an interior chart if $\varphi \left( U\right)$ is an open subset of ${\mathbb{R}}^{n}$ (which includes the case of an open subset of ${\mathbb{H}}^{n}$ that does not intersect $\partial {\mathbb{H}}^{n}$ ), and a boundary chart if $\varphi \left( U\right)$ is an open subset of ${\mathbb{H}}^{n}$ such that $\varphi \left( U\right)  \cap  \partial {\mathbb{H}}^{n} \neq  \varnothing$ . A boundary chart whose image is a set of the form ${B}_{r}\left( x\right)  \cap  {\mathbb{H}}^{n}$ for some $x \in  \partial {\mathbb{H}}^{n}$ and $r > 0$ is called a coordinate half-ball.

A point $p \in  M$ is called an interior point of $M$ if it is in the domain of some interior chart. It is a boundary point of $M$ if it is in the domain of a boundary chart that sends $p$ to $\partial {\mathbb{H}}^{n}$ . The boundary of $M$ (the set of all its boundary points) is denoted by $\partial M$ ; similarly, its interior, the set of all its interior points, is denoted by Int $M$ .

It follows from the definition that each point $p \in  M$ is either an interior point or a boundary point: if $p$ is not a boundary point, then either it is in the domain of an interior chart or it is in the domain of a boundary chart $\left( {U,\varphi }\right)$ such that $\varphi \left( p\right)  \notin  \partial {\mathbb{H}}^{n}$ , in which case the restriction of $\varphi$ to $U \cap  {\varphi }^{-1}\left( {\operatorname{Int}{\mathbb{H}}^{n}}\right)$ is an interior chart whose domain contains $p$ . However, it is not obvious that a given point cannot be simultaneously an interior point with respect to one chart and a boundary point with respect to another. In fact, this cannot happen, but the proof requires more machinery than we have available at this point. For convenience, we state the theorem here.
\end{definition}

\begin{theorem}[Topological Invariance of the Boundary]
\label{thm:1-37}
\lean{ModelWithCorners.disjoint_interior_boundary}
\mathlibok
\dcref{ch1:1.37}
If $M$ is a topological manifold with boundary, then each point of $M$ is either a boundary point or an interior point, but not both. Thus $\partial M$ and Int $M$ are disjoint sets whose union is $M$ .
\end{theorem}


\section{Smooth Structures on Manifolds with Boundary}

\begin{definition}
\label{def:1-6-extra-2}
\lean{SmoothManifoldWithBoundary}
\leanok
Now let $M$ be a topological manifold with boundary. As in the manifold case, a smooth structure for $M$ is defined to be a maximal smooth atlas-a collection of charts whose domains cover $M$ and whose transition maps (and their inverses) are smooth in the sense just described. With such a structure, $M$ is called a smooth manifold with boundary. Every smooth manifold is automatically a smooth manifold with boundary (whose boundary is empty).
\end{definition}

\begin{definition}
\label{def:1-6-extra-3}
\lean{IsRegularCoordinateHalfBall}
\leanok
Just as for smooth manifolds, if $M$ is a smooth manifold with boundary, any chart in the given smooth atlas is called a smooth chart for $M$ . Smooth coordinate balls, smooth coordinate half-balls, and regular coordinate balls in M are defined in the obvious ways. In addition, a subset $B \subseteq  M$ is called a regular coordinate half-ball if there is a smooth coordinate half-ball ${B}^{\prime } \supseteq  \bar{B}$ and a smooth coordinate map $\varphi  : {B}^{\prime } \rightarrow  {\mathbb{H}}^{n}$ such that for some ${r}^{\prime } > r > 0$ we have

\[
\varphi \left( B\right)  = {B}_{r}\left( 0\right)  \cap  {\mathbb{H}}^{n},\;\varphi \left( \bar{B}\right)  = {\bar{B}}_{r}\left( 0\right)  \cap  {\mathbb{H}}^{n},\;\text{ and }\;\varphi \left( {B}^{\prime }\right)  = {B}_{{r}^{\prime }}\left( 0\right)  \cap  {\mathbb{H}}^{n}.
\]
\end{definition}

\begin{exercise}
\label{exer:1-44}
\lean{manifoldInterior_boundaryless}
\leanok
\dcref{ch1:1.44}
Suppose $M$ is a smooth $n$ -manifold with boundary and $U$ is an open subset of $M$ . Prove the following statements:

(a) $U$ is a topological $n$ -manifold with boundary, and the atlas consisting of all smooth charts $\left( {V,\varphi }\right)$ for $M$ such that $V \subseteq  U$ defines a smooth structure on $U$ . With this topology and smooth structure, $U$ is called an open submanifold with boundary.

(b) If $U \subseteq  \operatorname{Int}M$ , then $U$ is actually a smooth manifold (without boundary); in this case we call it an open submanifold of $M$ .

(c) Int $M$ is an open submanifold of $M$ (without boundary).
\end{exercise}

\begin{proposition}
\label{prop:1-45}
\lean{ModelWithCorners.boundary_of_boundaryless_left}
\mathlibok
\dcref{ch1:1.45}
Suppose ${M}_{1},\ldots ,{M}_{k}$ are smooth manifolds and $N$ is a smooth manifold with boundary. Then ${M}_{1} \times  \cdots  \times  {M}_{k} \times  N$ is a smooth manifold with boundary, and $\partial \left( {{M}_{1} \times  \cdots  \times  {M}_{k} \times  N}\right)  = {M}_{1} \times  \cdots  \times  {M}_{k} \times  \partial N$ .
\end{proposition}

\begin{proof}
Problem 1-12.
\end{proof}

\begin{theorem}[Smooth Invariance of the Boundary]
\label{thm:1-46}
\lean{smooth_boundary_chart_frontier_independence}
\leanok
\dcref{ch1:1.46}
Suppose $M$ is a smooth manifold with boundary and $p \in  M$ . If there is some smooth chart $\left( {U,\varphi }\right)$ for $M$ such that $\varphi \left( U\right)  \subseteq  {\mathbb{H}}^{n}$ and $\varphi \left( p\right)  \in  \partial {\mathbb{H}}^{n}$ , then the same is true for every smooth chart whose domain contains $p$ .
\end{theorem}

\begin{proof}
\leanok
Suppose on the contrary that $p$ is in the domain of a smooth interior chart $\left( {U,\psi }\right)$ and also in the domain of a smooth boundary chart $\left( {V,\varphi }\right)$ such that $\varphi \left( p\right)  \in \; \partial {\mathbb{H}}^{n}$ . Let $\tau  = \varphi  \circ  {\psi }^{-1}$ denote the transition map; it is a homeomorphism from $\psi \left( {U \cap  V}\right)$ to $\varphi \left( {U \cap  V}\right)$ . The smooth compatibility of the charts ensures that both $\tau$ and ${\tau }^{-1}$ are smooth, in the sense that locally they can be extended, if necessary, to smooth maps defined on open subsets of ${\mathbb{R}}^{n}$ .

Write ${x}_{0} = \psi \left( p\right)$ and ${y}_{0} = \varphi \left( p\right)  = \tau \left( {x}_{0}\right)$ . There is some neighborhood $W$ of ${y}_{0}$ in ${\mathbb{R}}^{n}$ and a smooth function $\eta  : W \rightarrow  {\mathbb{R}}^{n}$ that agrees with ${\tau }^{-1}$ on $W \cap  \varphi \left( {U \cap  V}\right)$ . On the other hand, because we are assuming that $\psi$ is an interior chart, there is an open Euclidean ball $B$ that is centered at ${x}_{0}$ and contained in $\varphi \left( {U \cap  V}\right)$ , so $\tau$ itself is smooth on $B$ in the ordinary sense. After shrinking $B$ if necessary, we may assume that $B \subseteq  {\tau }^{-1}\left( W\right)$ . Then ${\left. \eta  \circ  \tau \right| }_{B} = {\left. {\tau }^{-1} \circ  \tau \right| }_{B} = {\operatorname{Id}}_{B}$ , so it follows from the chain rule that ${D\eta }\left( {\tau \left( x\right) }\right)  \circ  {D\tau }\left( x\right)$ is the identity map for each $x \in  B$ . Since ${D\tau }\left( x\right)$ is a square matrix, this implies that it is nonsingular. It follows from Corollary C. 36 that $\tau$ (considered as a map from $B$ to ${\mathbb{R}}^{n}$ ) is an open map, so $\tau \left( B\right)$ is an open subset of ${\mathbb{R}}^{n}$ that contains ${y}_{0} = \varphi \left( p\right)$ and is contained in $\varphi \left( V\right)$ . This contradicts the assumption that $\varphi \left( V\right)  \subseteq  {\mathbb{H}}^{n}$ and $\varphi \left( p\right)  \in  \partial {\mathbb{H}}^{n}$ .
\end{proof}


\section{Problems}

\begin{exercise}
\label{prob:1-1}
\lean{problem1LineWithTwoOrigins_origins_nhds_inter}
\leanok
1-1. Let $X$ be the set of all points $\left( {x, y}\right)  \in  {\mathbb{R}}^{2}$ such that $y =  \pm  1$ , and let $M$ be the quotient of $X$ by the equivalence relation generated by $\left( {x, - 1}\right)  \sim  \left( {x,1}\right)$ for all $x \neq  0$ . Show that $M$ is locally Euclidean and second-countable, but not Hausdorff. (This space is called the line with two origins.)
\end{exercise}

\begin{exercise}
\label{prob:1-2}
\lean{Sigma.t2Space}
\mathlibok
1-2. Show that a disjoint union of uncountably many copies of $\mathbb{R}$ is locally Euclidean and Hausdorff, but not second-countable.
\end{exercise}

\begin{exercise}
\label{prob:1-4}
\lean{finite_option_bounding_family}
\leanok
1-4. Let $M$ be a topological manifold, and let $\mathcal{U}$ be an open cover of $M$ .

(a) Assuming that each set in $\mathcal{U}$ intersects only finitely many others, show that $\mathcal{U}$ is locally finite.

(b) Give an example to show that the converse to (a) may be false.

(c) Now assume that the sets in $\mathcal{U}$ are precompact in $M$ , and prove the converse: if $\mathcal{U}$ is locally finite, then each set in $\mathcal{U}$ intersects only finitely many others.
\end{exercise}

\begin{exercise}
\label{prob:1-5}
\lean{isSigmaCompact_connectedComponent_of_paracompact_t2_euclidean}
\leanok
1-5. Suppose $M$ is a locally Euclidean Hausdorff space. Show that $M$ is second-countable if and only if it is paracompact and has countably many connected components. [Hint: assuming $M$ is paracompact, show that each component of $M$ has a locally finite cover by precompact coordinate domains, and extract from this a countable subcover.]
\end{exercise}

\begin{exercise}
\label{prob:1-6}
\lean{center_chart_transition_eqOn_unit_ball}
\leanok
1-6. Let $M$ be a nonempty topological manifold of dimension $n \geq  1$ . If $M$ has a smooth structure, show that it has uncountably many distinct ones. [Hint: first show that for any $s > 0,{F}_{s}\left( x\right)  = {\left| x\right| }^{s - 1}x$ defines a homeomorphism from ${\mathbb{B}}^{n}$ to itself, which is a diffeomorphism if and only if $s = 1$ .]
\end{exercise}

\begin{exercise}
\label{prob:1-7}
\lean{stereographic_two_chart_same_smooth_structure}
\leanok
1-7. Let $N$ denote the north pole $\left( {0,\ldots ,0,1}\right)  \in  {\mathbb{S}}^{n} \subseteq  {\mathbb{R}}^{n + 1}$ , and let $S$ denote the south pole $\left( {0,\ldots ,0, - 1}\right)$ . Define the stereographic projection $\sigma  : {\mathbb{S}}^{n} \smallsetminus  \{ N\}  \rightarrow  {\mathbb{R}}^{n}$ by

\[
\sigma \left( {{x}^{1},\ldots ,{x}^{n + 1}}\right)  = \frac{\left( {x}^{1},\ldots ,{x}^{n}\right) }{1 - {x}^{n + 1}}.
\]

Let $\widetilde{\sigma }\left( x\right)  =  - \sigma \left( {-x}\right)$ for $x \in  {\mathbb{S}}^{n} \smallsetminus  \{ S\}$ .

(a) For any $x \in  {\mathbb{S}}^{n} \smallsetminus  \{ N\}$ , show that $\sigma \left( x\right)  = u$ , where $\left( {u,0}\right)$ is the point where the line through $N$ and $x$ intersects the linear subspace where ${x}^{n + 1} = 0$ (Fig. 1.13). Similarly, show that $\widetilde{\sigma }\left( x\right)$ is the point where the line through $S$ and $x$ intersects the same subspace. (For this reason, $\widetilde{\sigma }$ is called stereographic projection from the south pole.)

(b) Show that $\sigma$ is bijective, and

\[
{\sigma }^{-1}\left( {{u}^{1},\ldots ,{u}^{n}}\right)  = \frac{\left( 2{u}^{1},\ldots ,2{u}^{n},{\left| u\right| }^{2} - 1\right) }{{\left| u\right| }^{2} + 1}.
\]

(c) Compute the transition map $\widetilde{\sigma } \circ  {\sigma }^{-1}$ and verify that the atlas consisting of the two charts $\left( {{\mathbb{S}}^{n}\smallsetminus \{ N\} ,\sigma }\right)$ and $\left( {{\mathbb{S}}^{n}\smallsetminus \{ S\} ,\widetilde{\sigma }}\right)$ defines a smooth structure on ${\mathbb{S}}^{n}$ . (The coordinates defined by $\sigma$ or $\widetilde{\sigma }$ are called stereographic coordinates.)

(d) Show that this smooth structure is the same as the one defined in Example 1.31.

(Used on pp. 201, 269, 301, 345, 347, 450.)

![45_428_187_679_414_0.jpg](images/45_428_187_679_414_0.jpg)

Fig. 1.13 Stereographic projection
\end{exercise}

\begin{exercise}
\label{prob:1-8}
\lean{IsAngleFunction}
\leanok
1-8. By identifying ${\mathbb{R}}^{2}$ with $\mathbb{C}$ , we can think of the unit circle ${\mathbb{S}}^{1}$ as a subset of the complex plane. An angle function on a subset $U \subseteq  {\mathbb{S}}^{1}$ is a continuous function $\theta  : U \rightarrow  \mathbb{R}$ such that ${e}^{{i\theta }\left( z\right) } = z$ for all $z \in  U$ . Show that there exists an angle function $\theta$ on an open subset $U \subseteq  {\mathbb{S}}^{1}$ if and only if $U \neq  {\mathbb{S}}^{1}$ . For any such angle function, show that $\left( {U,\theta }\right)$ is a smooth coordinate chart for ${\mathbb{S}}^{1}$ with its standard smooth structure. (Used on pp. 37,152,176.)
\end{exercise}

\begin{exercise}
\label{prob:1-9}
\lean{complexProjectiveSpaceIsManifold}
\leanok
1-9. Complex projective n-space, denoted by ${\mathbb{{CP}}}^{n}$ , is the set of all 1-dimensional complex-linear subspaces of ${\mathbb{C}}^{n + 1}$ , with the quotient topology inherited from the natural projection $\pi  : {\mathbb{C}}^{n + 1} \smallsetminus  \{ 0\}  \rightarrow  {\mathbb{{CP}}}^{n}$ . Show that ${\mathbb{{CP}}}^{n}$ is a compact ${2n}$ -dimensional topological manifold, and show how to give it a smooth structure analogous to the one we constructed for ${\mathbb{{RP}}}^{n}$ . (We use the correspondence

\[
\left( {{x}^{1} + i{y}^{1},\ldots ,{x}^{n + 1} + i{y}^{n + 1}}\right)  \leftrightarrow  \left( {{x}^{1},{y}^{1},\ldots ,{x}^{n + 1},{y}^{n + 1}}\right)
\]

to identify ${\mathbb{C}}^{n + 1}$ with ${\mathbb{R}}^{{2n} + 2}$ .) (Used on pp. 48,96,172,560,561.)
\end{exercise}

\begin{exercise}
\label{prob:1-10}
\lean{standard_chart_matrix_unique}
\leanok
\uses{ex:1-36}
1-10. Let $k$ and $n$ be integers satisfying $0 < k < n$ , and let $P, Q \subseteq  {\mathbb{R}}^{n}$ be the linear subspaces spanned by $\left( {{e}_{1},\ldots ,{e}_{k}}\right)$ and $\left( {{e}_{k + 1},\ldots ,{e}_{n}}\right)$ , respectively, where ${e}_{i}$ is the $i$ th standard basis vector for ${\mathbb{R}}^{n}$ . For any $k$ -dimensional subspace $S \subseteq  {\mathbb{R}}^{n}$ that has trivial intersection with $Q$ , show that the coordinate representation $\varphi \left( S\right)$ constructed in Example 1.36 is the unique $\left( {n - k}\right)  \times  k$ matrix $B$ such that $S$ is spanned by the columns of the matrix $\left( \begin{matrix} {I}_{k} \\  B \end{matrix}\right)$ , where ${I}_{k}$ denotes the $k \times  k$ identity matrix.
\end{exercise}

\begin{exercise}
\label{prob:1-11}
\lean{closed_unit_ball_subtype_val_contMDiff}
\leanok
\uses{prob:1-7}
1-11. Let $M = {\overline{\mathbb{B}}}^{n}$ , the closed unit ball in ${\mathbb{R}}^{n}$ . Show that $M$ is a topological manifold with boundary in which each point in ${\mathbb{S}}^{n - 1}$ is a boundary point and each point in ${\mathbb{B}}^{n}$ is an interior point. Show how to give it a smooth structure such that every smooth interior chart is a smooth chart for the standard smooth structure on ${\mathbb{B}}^{n}$ . [Hint: consider the map $\pi  \circ  {\sigma }^{-1} : {\mathbb{R}}^{n} \rightarrow  {\mathbb{R}}^{n}$ , where $\sigma  : {\mathbb{S}}^{n} \rightarrow  {\mathbb{R}}^{n}$ is the stereographic projection (Problem 1-7) and $\pi$ is a projection from ${\mathbb{R}}^{n + 1}$ to ${\mathbb{R}}^{n}$ that omits some coordinate other than the last.]
\end{exercise}

\begin{exercise}
\label{prob:1-12}
\lean{prod_boundaryless_left_isManifold_with_boundary}
\leanok
\uses{prop:1-45}
1-12. Prove Proposition 1.45 (a product of smooth manifolds together with one smooth manifold with boundary is a smooth manifold with boundary).
\end{exercise}

